\section*{Part VII: The Bayes Institute Distribution System}

\subsection*{56. The Starting Constraint Is Attention, Not Merely Capital}

Bayes Institute commences operations as a bootstrapped, asset-light technical education and research studio. Its founder maintains a demanding full-time institutional role, dedicating approximately 15 to 25 sustainable weekly hours to Bayes across curriculum development, teaching, admissions vetting, and institutional operations. Under these operating conditions, \textbf{founder attention is the binding organizational constraint}, far more restrictive than financial capital.

A distribution channel that demands daily short-form social content creation, unrestricted 30-minute discovery calls with casual inquiries, or extensive offline conference travel is economically ruinous, even if its direct monetary advertising cost is zero. The distribution objective is to architect an acquisition mechanism that repeatedly encounters a narrowly defined buying situation, converts eligible candidates with high conversion efficiency, and preserves the founder's scarce hours for uncompromised instructional delivery. While the institute's intellectual identity and curriculum standards are globally legible, its initial customer acquisition is tightly focused geographically and behaviorally.

\subsection*{57. Mapping the 3 Hypothetical Customer Routes}

Bayes Institute addresses 3 distinct customer systems, each governed by unique purchasing triggers, evaluation cycles, and unit economics:
\begin{enumerate}[leftmargin=*,itemsep=1pt]
  \item \textbf{Direct-to-Consumer (D2C) Technical Professionals:}
  \begin{itemize}[leftmargin=*,itemsep=1pt]
    \item \textit{Operational trigger:} An experienced backend or machine-learning engineer assumes responsibility for a customer-facing LLM application and discovers that generic tutorials provide zero guidance for debugging silent evaluation regressions in production.
    \item \textit{Access pathways:} Deep technical failure teardowns, specialized developer forums, engineering peer recommendations, and high-intent long-tail search.
    \item \textit{Decision cycle:} 7 to 21 days; self-funded or reimbursed via personal professional-development budgets.
  \end{itemize}
  \item \textbf{Engineering College Partnerships:}
  \begin{itemize}[leftmargin=*,itemsep=1pt]
    \item \textit{Operational trigger:} Academic leadership (Dean of Engineering or Department Chair) recognizes that existing computer-science curricula lack rigorous, production-grade AI systems electives and seeks an accredited industry lab module.
    \item \textit{Access pathways:} Academic alumni introductions, faculty seminars, and structured institutional pilot proposals.
    \item \textit{Decision cycle:} 6 to 12 months; governed by institutional board approvals, academic semester calendars, and student laboratory fee structures.
  \end{itemize}
  \item \textbf{Enterprise Engineering Teams:}
  \begin{itemize}[leftmargin=*,itemsep=1pt]
    \item \textit{Operational trigger:} An engineering director or VP faces severe reliability bottlenecks, customer complaints, or compliance review failures stemming from uncontrolled model hallucinations.
    \item \textit{Access pathways:} Account-based technical outreach, practitioner referrals, and bounded architectural diagnostic reports.
    \item \textit{Decision cycle:} 30 to 90 days; funded via departmental engineering budgets or vendor advisory retainers.
  \end{itemize}
\end{enumerate}

These 3 customer pathways cannot be evaluated on a single conversion dashboard. A college partnership meeting converts at a vastly lower velocity than a D2C inquiry, but can contract for 60 learner seats simultaneously. The initial distribution priority is to validate the D2C motion decisively while conducting patient discovery across institutional routes.

\subsection*{58. The 1st Canonical Evidence Asset}

To establish technical authority without commercial shouting, Bayes Institute's initial distribution asset is an open, deeply rigorous engineering case study: \textit{``Why a High-Scoring Retrieval Demo Fails in Production: And How to Design an Evaluation Harness That Catches It.''}

This canonical asset includes:
\begin{itemize}[leftmargin=*,itemsep=1pt]
  \item Fully runnable benchmark code and synthetic retrieval datasets.
  \item Documented edge-case failure traces (such as semantic drift and document chunking boundary errors).
  \item An auditable evaluation rubric comparing heuristic, cross-encoder, and model-based judges.
  \item A concise 12-minute technical walk-through analyzing mathematical trade-offs between precision and latency.
  \item A downloadable diagnostic self-assessment that allows an engineer to score their team's current evaluation rigor.
\end{itemize}

This asset achieves extraordinary operational leverage. It simultaneously establishes academic credibility (Chapter 19), attracts organic practitioner search demand (Chapter 20), anchors consultative admissions vetting (Chapter 21), and serves as Laboratory Module 1 in the active curriculum (Chapter 12). However, a single asset does not constitute a distribution engine until a repeatable route delivers it to qualified practitioners.

\subsection*{59. A Disciplined 30-Day D2C Pilot Experiment}

Rather than testing 6 channels superficially, the founder executes a controlled 30-day distribution experiment:
\begin{itemize}[leftmargin=*,itemsep=1pt]
  \item \textbf{Target profile:} Senior backend and data engineers actively maintaining retrieval-augmented generation pipelines in mid-sized technology firms.
  \item \textbf{Core offer:} A 2-weekend, 12-hour advanced practicum: \textit{Production RAG Evaluation and Reliability}.
  \item \textbf{Execution motion:} The founder shares the canonical technical study with 10 senior engineering peers for candid critique; contributes the code repository to 1 specialized practitioner forum that explicitly permits open-source sharing; and offers 8 paid seats to qualified applicants.
  \item \textbf{Measurement:} The founder logs every hour expended, every inbound inquiry, all technical objections, payment confirmations, and required onboarding support.
\end{itemize}

Success is defined strictly by whether 8 fully qualified engineers pay the tuition fee, complete the practicum, and verify that the learning transformed their production systems. If zero prospects enroll, the experiment succeeds informatively: allowing the founder to isolate whether the failure was driven by pricing, timing, prerequisite miscalibration, or inadequate distribution reach.

\subsection*{60. Illustrative D2C Unit Distribution Economics}

Consider the reconciled unit economics of an 8-seat pilot cohort:
\begin{align*}
  \text{Gross tuition (8 seats at ₹20,000)} &= \text{₹}1,60,000 \\
  \text{Direct delivery and labs} &= \text{₹}60,000 \\
  \hline
  \text{Pre-acquisition contribution} &= \text{₹}1,00,000 \\
  \text{Direct cash acquisition spend} &= \text{₹}12,000 \\
  \text{Founder time (26 hrs at ₹1,000/hr)} &= \text{₹}26,000 \\
  \hline
  \text{Economic acquisition cost} &= \text{₹}38,000 \\
  \text{Post-acquisition contribution} &= \text{₹}62,000
\end{align*}

The economic acquisition cost is ₹4,750 per learner (₹1,500 cash CAC plus ₹3,250 founder-labor CAC). The cohort yields ₹62,000 in net cash contribution to cover technological infrastructure and reserve funds.

Now examine the marginal sensitivity. If Cohort 1 relies heavily on warm professional networks, Cohort 2 may require colder outreach or paid channels. If acquiring Cohort 2 demands ₹60,000 in cash advertising spend and the same 26 founder hours, total acquisition costs rise to ₹86,000: collapsing post-acquisition contribution from ₹62,000 to ₹14,000. Scaling cannot be justified until the venture confirms whether subsequent cohorts can be acquired at sustainable marginal economics.

\subsection*{61. The College Channel Begins with Authorized Discovery}

While the founder possesses previous experience with institutional university placement programs and curriculum advisory projects, those exploratory interactions do not constitute a repeatable distribution channel. Institutional academic sales follow bureaucratic protocols and multi-year calendar rhythms.

A disciplined college acquisition test proceeds as follows:
\begin{itemize}[leftmargin=*,itemsep=1pt]
  \item Identify 3 autonomous engineering colleges with established industry-sponsored laboratories.
  \item Engage the Dean of Academic Affairs and Department Chair around a specific operational problem: bridging the gap between theoretical machine learning and enterprise production standards.
  \item Propose an authorized, non-credit weekend laboratory module (e.g. 30 student seats funded via an existing laboratory development budget).
  \item Establish explicit contractual boundaries: student eligibility prerequisites, grading rubrics, faculty co-teaching commitments, and student data privacy compliance.
\end{itemize}

The metric of success is an authorized institutional agreement with confirmed budget allocation. Slower than D2C, a successful college channel delivers aggregated blocks of learners with zero marginal consumer ad spend, diversifying institutional revenue.

\subsection*{62. Enterprise Distribution Begins with an Auditable Diagnostic}

Marketing an undifferentiated ``Corporate AI Training'' seminar pits Bayes Institute against large, established corporate training vendors and generates minimal urgency. Enterprise engineering leaders do not purchase broad training; they contract solutions to acute operational crises.

The enterprise go-to-market motion initiates with an auditable technical assessment:
\begin{itemize}[leftmargin=*,itemsep=1pt]
  \item Identify 20 engineering organizations deploying mission-critical LLM features.
  \item Provide the VP of Engineering with a concise, 3-page technical audit rubric: \textit{``The Enterprise Evaluation Readiness Checklist.''}
  \item Offer a bounded, paid 2-day technical diagnostic: auditing the organization's existing evaluation pipelines, benchmarking hallucination rates, and delivering an executive remediation memo.
  \item Use the diagnostic findings to scope a dedicated 4-week team practicum for their senior engineers.
\end{itemize}

This diagnostic-led motion positions the institute as an elite technical advisory rather than an outsourced training vendor, commands premium pricing (₹3,00,000 to ₹5,00,000 per engagement), and establishes a defensible enterprise beachhead.
